\section{Planning：长程任务中的策略控制器}

\subsection{规划的作用边界}
课程将 Planning 定义为 ``在约束下分配未来动作预算''，而非生成静态待办清单。高质量规划必须允许动态重规划（dynamic replanning），因为环境反馈会持续改变最优动作序列。

\begin{figure}[H]
\centering
\includegraphics[width=0.84\textwidth]{slide-05-planning-loop.jpg}
\caption{视频约 01:08:20 讲义页：规划需要与执行和反馈形成闭环，而非一次性展开}
\end{figure}

\begin{importantbox}{规划闭环的三个必要环节}
\begin{enumerate}
    \item 目标分解：将最终目标映射为可验证中间里程碑；
    \item 执行监控：持续比较预期状态与真实状态；
    \item 计划改写：在偏差超阈值时触发重规划。
\end{enumerate}
\end{importantbox}

\subsection{规划与搜索的结合}
Yu Su 的观点是 ``规划不是搜索的替代，搜索也不是规划的替代''。规划决定搜索空间结构，搜索决定局部动作最优。系统工程中应按任务复杂度选择是否引入 MCTS、beam search 或启发式候选生成。

\begin{knowledgebox}{何时需要显式搜索}
\begin{itemize}
    \item 行动后果延迟且不可逆；
    \item 局部最优陷阱明显；
    \item 单步评分器信号弱、噪声大。
\end{itemize}
\end{knowledgebox}

\begin{warningbox}{过度规划的代价}
过细粒度规划会占用大量上下文与计算预算，并削弱执行阶段的实时适应能力。规划粒度应与环境变化速率匹配。
\end{warningbox}

\subsection{规划器伪代码}
\begin{lstlisting}[caption={一个简化的重规划循环，强调“偏差触发式更新”}]
state = observe()
plan = planner.make_plan(goal, state)

for step in range(budget):
    action = planner.next_action(plan, state)
    result = env.execute(action)
    state = update_state(state, result)
    memory.write(result)

    if planner.need_replan(state, goal):
        plan = planner.replan(goal, state, memory.retrieve(goal))
\end{lstlisting}

\subsection{本章小结}
Planning 的关键不在于 ``提前想得多完整''，而在于 ``执行中能否快速修正''。动态重规划是长程 Agent 成败分水岭。

